\documentclass[12pt, a4paper]{report}
\usepackage[utf8]{inputenc}
\usepackage[T1]{fontenc}
\usepackage[brazil]{babel}
\usepackage{graphicx}
\usepackage{geometry}
\usepackage{titlesec}
\usepackage{hyperref}
\usepackage{booktabs}
\usepackage{amsmath}
\usepackage{float}
\usepackage{natbib}
\usepackage{xcolor}
\usepackage{longtable}
\usepackage{setspace} 

% Configuração das margens (ABNT)
\geometry{top=3cm, bottom=2cm, left=3cm, right=2cm}
\onehalfspacing

% Configuração de Links
\hypersetup{
    colorlinks=true,
    linkcolor=blue,
    filecolor=magenta,      
    urlcolor=cyan,
    citecolor=blue,
}

% Título e Autor
\title{\textbf{Modelagem Preditiva e Auditoria da Qualidade de Dados da Hanseníase no Brasil: Uma Abordagem Comparativa entre Ensemble Learning e Deep Learning (2001-2025)}}
\author{Seu Nome}
\date{\today}

\begin{document}

\maketitle
\tableofcontents
\newpage

% ==================================================================
% CAPÍTULO 1: INTRODUÇÃO
% ==================================================================
\chapter{Introdução}

\section{Contextualização e Problema de Pesquisa}
A hanseníase é uma doença infectocontagiosa crônica de evolução lenta, causada pelo bacilo \textit{Mycobacterium leprae}, que representa um desafio persistente e complexo para a saúde pública no Brasil. O país ocupa a segunda posição mundial em número de novos casos detectados anualmente, ficando atrás apenas da Índia, e concentra mais de 90\% dos casos registrados nas Américas \citep{who2021zero}. A principal complicação da doença é o desenvolvimento de incapacidades físicas, decorrentes do dano neural periférico (neurite), que podem evoluir para deformidades irreversíveis se não diagnosticadas e tratadas precocemente. Essas sequelas não apenas impactam a funcionalidade do indivíduo, mas também perpetuam o estigma histórico e a exclusão social associados à enfermidade \citep{brasil2022pcdt}.

O Sistema de Informação de Agravos de Notificação (SINAN) constitui a espinha dorsal da vigilância epidemiológica no país, coletando dados compulsórios sobre cada caso diagnosticado. Contudo, a qualidade desses registros administrativos é historicamente heterogênea, apresentando desafios como subnotificação, duplicidades e incompletude de campos essenciais para a avaliação de incapacidade. A pandemia de COVID-19 (2020-2023) impôs um choque exógeno sem precedentes a este sistema. A reorganização dos serviços de saúde para o combate ao vírus resultou na descontinuidade de ações programáticas, como a busca ativa de casos e o acompanhamento de contatos, gerando um represamento de diagnósticos. Estima-se que milhares de casos deixaram de ser notificados ou foram notificados com avaliações físicas incompletas durante este período, criando um "apagão" de informações sobre a real magnitude das incapacidades instaladas na população \citep{saldanha2021ciencia, tang2021bigdata}.

\section{Justificativa}
A complexidade dos determinantes da hanseníase — que envolvem fatores biológicos, clínicos e socioeconômicos — e o volume massivo de dados acumulados no SINAN ao longo de duas décadas configuram um cenário de \textit{Big Data} em saúde. As metodologias estatísticas tradicionais, muitas vezes limitadas a modelos lineares e análises descritivas, podem falhar em capturar padrões não-lineares sutis e interações complexas entre variáveis. A aplicação de técnicas avançadas de Ciência de Dados (\textit{Data Science}) e Aprendizado de Máquina (\textit{Machine Learning}) oferece uma nova perspectiva analítica. Algoritmos capazes de aprender com dados históricos complexos podem ser utilizados não apenas para predição de risco individual, mas também como ferramentas de auditoria \citep{goldstein2020convergence}. Ao modelar o comportamento "esperado" da doença com base em décadas de dados, é possível identificar desvios anômalos no padrão de notificações recentes, quantificando a subnotificação qualificada (erro de classificação) e fornecendo evidências para direcionar políticas de busca ativa e recuperação de dados.

\section{Objetivos}
O objetivo geral deste trabalho é desenvolver um \textit{framework} robusto de modelagem preditiva para estimar o risco de Incapacidade Física (GIF) na alta de pacientes com hanseníase, utilizando-o posteriormente para avaliar a integridade e a qualidade dos dados no período pós-pandemia.

\textbf{Objetivos Específicos:}
\begin{enumerate}
    \item Estruturar um \textit{pipeline} reprodutível de engenharia de dados para o tratamento, limpeza e normalização de notificações do SINAN abrangendo o período de 2001 a 2025.
    \item Implementar e avaliar técnicas de seleção de atributos, especificamente a Eliminação Recursiva de Features (RFE), para identificar os marcadores clínicos e demográficos mais robustos de prognóstico de incapacidade.
    \item Comparar a eficácia e a eficiência computacional de algoritmos de \textit{Ensemble Learning} (\textit{Bagging} com Random Forest e \textit{Boosting} com LightGBM) versus abordagens de \textit{Deep Learning} (\textit{Gated Recurrent Units} - GRU) em dados tabulares epidemiológicos desbalanceados.
    \item Quantificar o "apagão de dados" de incapacidade física grave no período de 2020 a 2025 através da análise de resíduos do modelo, estimando o volume de casos com grau de incapacidade subnotificado.
\end{enumerate}

% ==================================================================
% CAPÍTULO 2: FUNDAMENTAÇÃO TEÓRICA
% ==================================================================
\chapter{Fundamentação Teórica e Estado da Arte}

\section{A Hanseníase e seus Indicadores de Gravidade}
A hanseníase manifesta-se através de lesões dermatológicas e neurológicas. Para fins operacionais de tratamento (poliquimioterapia), a doença é classificada em **Paucibacilar (PB)** e **Multibacilar (MB)**, sendo esta última associada a maior carga bacilar, maior risco de transmissão e maior probabilidade de reações inflamatórias que causam dano neural.

O **Grau de Incapacidade Física (GIF)** é o padrão-ouro internacional para avaliação da severidade e monitoramento da doença. Ele é avaliado através de exame físico simplificado (avaliação neurológica) nos olhos, mãos e pés, e classificado em uma escala de 0 a 2:
\begin{itemize}
    \item \textbf{Grau 0:} Nenhuma perda de sensibilidade ou força muscular.
    \item \textbf{Grau 1:} Diminuição ou perda de sensibilidade protetora nas mãos ou pés, sem deformidades visíveis.
    \item \textbf{Grau 2:} Presença de deformidades visíveis (como mãos em garra, lagoftalmo, reabsorção óssea) ou alterações visuais graves \citep{brasil2017guia}.
\end{itemize}
A evolução do GIF entre o momento do diagnóstico e a alta do tratamento é um indicador crítico da qualidade da assistência. O agravamento do grau (por exemplo, de 0 para 1 ou 2) durante o tratamento indica falha na prevenção de incapacidades (neurites não tratadas). A presença de Grau 2 já no diagnóstico é um indicador sentinela de detecção tardia, sugerindo que a doença evoluiu por anos sem intervenção. A análise da **matriz de transição** realizada preliminarmente neste estudo confirmou uma forte **inércia clínica**: pacientes que iniciam o tratamento com Grau 2 raramente revertem o quadro totalmente para Grau 0, o que reforça a necessidade absoluta do diagnóstico precoce para prevenção secundária.

\section{Algoritmos de Aprendizado de Máquina em Saúde}

\subsection{Ensemble Learning: Bagging e Boosting}
Métodos de \textit{Ensemble} (conjunto) baseiam-se na premissa de que a combinação de múltiplos modelos de aprendizado (modelos fracos) resulta em um modelo final com maior precisão e generalização do que qualquer um dos modelos individuais.

\subsubsection{Random Forest (Bagging)}
O Random Forest utiliza a técnica de \textit{Bagging} (Bootstrap Aggregating). Ele constrói múltiplas árvores de decisão independentes durante o treinamento. Cada árvore é treinada em um subconjunto aleatório dos dados (amostragem com reposição) e, em cada divisão de nó, considera apenas um subconjunto aleatório de \textit{features}. A predição final é obtida pela média (regressão) ou voto majoritário (classificação) das árvores.
\begin{itemize}
    \item \textbf{Vantagens:} É altamente robusto contra \textit{overfitting}, lida bem com dados não-lineares e interações entre variáveis sem necessidade de especificação prévia, e fornece medidas úteis de importância das variáveis. Por essas características, foi selecionado como o modelo \textit{baseline} deste estudo \citep{breiman2001random}.
\end{itemize}

\subsubsection{LightGBM (Gradient Boosting)}
O LightGBM (\textit{Light Gradient Boosting Machine}) pertence à família de algoritmos de \textit{Boosting}. Diferente do Bagging, o Boosting constrói o modelo de forma sequencial: cada nova árvore é treinada para corrigir os erros (resíduos) cometidos pelas árvores anteriores. O LightGBM introduz inovações algorítmicas importantes:
\begin{itemize}
    \item \textbf{Crescimento Leaf-wise (Por Folha):} A maioria dos algoritmos de árvore cresce nível por nível (\textit{level-wise}). O LightGBM escolhe a folha que maximiza a redução da função de perda para dividir. Isso permite que ele modele relações muito mais complexas e profundas, resultando em maior acurácia, embora com maior risco de \textit{overfitting} em pequenos datasets (o que é mitigado pelo grande volume de dados do SINAN).
    \item \textbf{Histogram-based Algorithm:} Agrupa valores contínuos de \textit{features} em bins discretos, acelerando drasticamente o treinamento e reduzindo o uso de memória.
    \item \textbf{Exclusive Feature Bundling (EFB):} Combina variáveis mutuamente exclusivas (que raramente são não-nulas simultaneamente) para reduzir a dimensionalidade sem perda de informação.
\end{itemize}
Essas características tornam o LightGBM o estado da arte para dados tabulares estruturados em larga escala \citep{ke2017lightgbm}.

\subsection{Deep Learning: Redes Neurais Recorrentes (GRU)}

As Redes Neurais Recorrentes (RNNs) são especializadas em processar dados sequenciais, mantendo um "estado oculto" (memória) que carrega informações de passos anteriores. No entanto, RNNs tradicionais sofrem com o problema do "gradiente evanescente", onde a rede esquece informações de longo prazo.

A **Gated Recurrent Unit (GRU)** é uma arquitetura otimizada que resolve esse problema utilizando mecanismos de controle de fluxo de informação chamados "portões":
\begin{itemize}
    \item \textbf{Update Gate ($z_t$):} Decide quanto da informação passada (memória anterior) deve ser mantida e quanto da nova informação deve ser adicionada.
    \item \textbf{Reset Gate ($r_t$):} Decide quanto da informação passada deve ser esquecida/ignorada.
\end{itemize}
Embora as GRUs sejam tipicamente aplicadas em séries temporais (como sinais vitais ou texto), sua aplicação em dados tabulares estáticos tem sido investigada. Neste contexto, o vetor de \textit{features} de um paciente é tratado como uma sequência curta ou como um passo temporal único com múltiplos canais. A hipótese é que a capacidade da GRU de realizar transformações não-lineares complexas através de seus portões poderia capturar interações sutis entre variáveis clínicas que modelos baseados em árvores poderiam perder \citep{cho2014gru, shahidi2024exploring}.

% ==================================================================
% CAPÍTULO 3: METODOLOGIA
% ==================================================================
\chapter{Metodologia}

\section{Fonte de Dados e Pré-processamento}
Os dados foram extraídos da base nacional do SINAN, compreendendo notificações de hanseníase de todo o território brasileiro no período de 2001 a 2025. O volume inicial de cerca de 984 mil registros passou por um rigoroso processo de limpeza e normalização:
\begin{enumerate}
    \item \textbf{Limpeza:} Remoção de registros duplicados (flag \texttt{NDUPLIC\_N}), casos transferidos e inconsistências de datas.
    \item \textbf{Definição do Target:} A variável alvo foi definida como o **GIF na Alta** (\texttt{AVAL\_ATU\_N\_CAT}). Para garantir a confiabilidade do treinamento supervisionado, registros com classificação "Ignorado" (9) ou "Não Avaliado" (3) foram excluídos da base de modelagem, resultando em um dataset final de aproximadamente 540 mil amostras rotuladas (Grau 0, 1 ou 2).
    \item \textbf{Engenharia de Features:} Criação de variáveis derivadas, como a duração do tratamento em dias (\texttt{DURACAO\_TRAT} = Data Alta - Data Diagnóstico), e normalização de variáveis. Variáveis categóricas (Sexo, Raça, Forma Clínica, etc.) foram convertidas para representação numérica via \textit{One-Hot Encoding} (para GRU) ou mantidas como categorias (para LightGBM, que lida nativamente com elas).
\end{enumerate}

\section{Seleção de Variáveis: Recursive Feature Elimination (RFE)}
Dados de saúde frequentemente contêm variáveis redundantes ou irrelevantes que podem introduzir ruído e aumentar a complexidade computacional. Para otimizar o modelo, utilizou-se o método de **Eliminação Recursiva de Features (RFE)**.
\begin{itemize}
    \item \textbf{Fundamentação Técnica:} O RFE é um método de seleção do tipo *wrapper*. Ele funciona treinando o modelo (neste caso, um Random Forest) com todas as variáveis disponíveis e calculando a importância de cada uma (baseada na impureza de Gini). A variável (ou conjunto de variáveis) menos importante é removida, e o modelo é retreinado com as restantes. Esse processo é repetido recursivamente até que o número desejado de variáveis seja atingido. Isso garante que o subconjunto final de \textit{features} seja aquele que maximiza o poder preditivo conjunto, considerando as interações entre elas \citep{guyon2002gene}.
    \item \textbf{Seleção Final:} O algoritmo convergiu para um subconjunto ótimo de 10 variáveis, descritas na Tabela \ref{tab:features}. A seleção confirmou a primazia de dados clínicos sobre demográficos.
\end{itemize}

\begin{table}[H]
\centering
\caption{Variáveis Selecionadas via RFE para Modelagem}
\label{tab:features}
\resizebox{\textwidth}{!}{%
\begin{tabular}{lll}
\toprule
\textbf{Variável} & \textbf{Tipo} & \textbf{Relevância Clínica e Justificativa} \\
\midrule
\texttt{AVALIA\_N\_CAT} & Categórica & GIF no diagnóstico: O preditor mais forte devido à inércia clínica. \\
\texttt{CLASSATUAL\_CAT} & Categórica & Define a carga bacilar (MB/PB) e o potencial inflamatório sistêmico. \\
\texttt{NERVOSAFET\_CAT} & Numérica & Quantificação direta do dano neural periférico (espessamento). \\
\texttt{BACILOSCOP\_CAT} & Categórica & Positividade correlaciona-se com formas Virchowianas e maior risco de dano. \\
\texttt{EPIS\_RACIO\_CAT} & Categórica & Reações (Tipo 1/2) são a principal causa aguda de dano neural durante o tratamento. \\
\texttt{IDADE\_ANOS} & Numérica & Fator de risco acumulativo; idosos têm pior recuperação neural. \\
\texttt{DURACAO\_TRAT} & Numérica & Tratamentos prolongados indicam complicações, interrupções ou falha terapêutica. \\
\texttt{DOSE\_RECEB\_CAT} & Categórica & Indicador de adesão e regularidade do tratamento. \\
\texttt{CS\_RACA\_CAT} & Categórica & Marcador proxy de determinantes sociais de saúde e vulnerabilidade. \\
\texttt{FORMACLINI\_CAT} & Categórica & Classificação clínica detalhada (Madrid) que reflete a resposta imune do hospedeiro. \\
\bottomrule
\end{tabular}
}
\end{table}

\section{Estratégia de Validação: Walk-Forward Validation}
A validação cruzada aleatória tradicional (\textit{k-fold}) assume que os dados são independentes e identicamente distribuídos (i.i.d). Em dados epidemiológicos, essa premissa é falha devido à autocorrelação temporal e mudanças de cenário (ex: mudanças em protocolos de saúde). O uso de \textit{k-fold} aleatório introduziria um viés de *data leakage* (vazamento de dados), onde o modelo aprenderia com dados do futuro para prever o passado.

Para garantir a validade metodológica, adotou-se a estratégia de **Validação Walk-Forward** (ou *Time Series Split*):
\begin{enumerate}
    \item \textbf{Treino e Otimização (2001-2019):} Utilizou-se janelas deslizantes nos dados históricos para a otimização de hiperparâmetros. O modelo é treinado em uma janela temporal $t$ e validado na janela $t+1$ (ex: Treino 2001-2005 -> Teste 2006; Treino 2001-2006 -> Teste 2007). Isso força o modelo a aprender a dinâmica temporal e evolutiva da endemia.
    \item \textbf{Teste de Estresse (2020-2025):} O conjunto de dados pós-pandemia foi mantido como um \textit{hold-out set} (totalmente isolado do processo de treino). O modelo final otimizado foi aplicado a este período para avaliar sua robustez frente ao \textit{concept drift} (mudança na distribuição dos dados) causado pelo choque da pandemia no sistema de saúde.
\end{enumerate}

% ==================================================================
% CAPÍTULO 4: RESULTADOS E DISCUSSÃO
% ==================================================================
\chapter{Resultados Experimentais e Discussão}



A avaliação dos modelos focou na capacidade de generalização no período de teste (2020-2025), um cenário marcado por instabilidade nos dados e possíveis alterações nos processos de notificação.

\section{Análise Comparativa de Algoritmos}
A Tabela \ref{tab:resultados} resume as métricas de performance no conjunto de teste. O foco principal foi equilibrar a Acurácia Global com o \textit{Recall} (Sensibilidade) para a classe mais grave (Grau 2), dado que o custo de não detectar um paciente com risco de deformidade é alto.

\begin{table}[H]
\centering
\caption{Performance dos Modelos no Conjunto de Teste (Pós-Pandemia)}
\label{tab:resultados}
\begin{tabular}{lcccc}
\toprule
\textbf{Modelo} & \textbf{Acurácia} & \textbf{Precisão (G2)} & \textbf{Recall (G2)} & \textbf{F1-Score (G2)} \\
\midrule
Random Forest & 71.16\% & 0.44 & 0.72 & 0.55 \\
Deep Learning (GRU) & 71.44\% & 0.42 & \textbf{0.73} & 0.53 \\
\textbf{LightGBM (Otimizado)} & \textbf{71.10\%} & \textbf{0.45} & 0.71 & \textbf{0.55} \\
\bottomrule
\end{tabular}
\end{table}

\subsection{Interpretação Técnica e Escolha do Modelo}
\begin{itemize}
    \item **Convergência de Performance:** Observou-se uma convergência na acurácia global (~71\%) entre as três abordagens distintas (Bagging, Boosting e Redes Neurais). Isso sugere fortemente que o limite de preditibilidade reside na qualidade intrínseca dos dados (ruído aleatório, erros de preenchimento humano, variabilidade biológica não capturada) e não na capacidade de modelagem dos algoritmos.
    \item **LightGBM vs. Random Forest:** O LightGBM apresentou um desempenho ligeiramente superior ao Random Forest em termos de precisão para a classe crítica (Grau 2), mantendo um alto *recall*. Sua principal vantagem, contudo, foi a eficiência computacional: o treinamento e a inferência foram significativamente mais rápidos, o que o torna mais escalável para aplicações em tempo real no SUS. A estratégia de crescimento *leaf-wise* permitiu um ajuste mais fino às fronteiras de decisão complexas das classes minoritárias.
    \item **Deep Learning (GRU):** A rede GRU obteve o maior *Recall* (0.73), demonstrando a mais alta sensibilidade para detectar casos graves. No entanto, sua precisão foi a menor (0.42), indicando uma taxa elevada de falsos positivos (classificar casos leves como graves). Este resultado corrobora a literatura recente que sugere que, para dados tabulares com muitas variáveis categóricas e sem estrutura sequencial longa, a complexidade computacional das redes neurais não se traduz necessariamente em ganho de acurácia sobre métodos de *Gradient Boosting* bem otimizados \citep{shwartz2022tabular}.
    \item **Decisão:** O **LightGBM** foi selecionado como o modelo final para a etapa de auditoria devido ao seu melhor equilíbrio entre sensibilidade e precisão (F1-Score), robustez e interpretabilidade superior em relação à "caixa preta" das redes neurais.
\end{itemize}



\section{Diagnóstico Epidemiológico: O "Apagão" de Dados}
A aplicação do modelo LightGBM aos dados de 2020-2025 revelou uma discrepância estatística significativa e preocupante na análise de resíduos (diferença entre o valor predito pelo modelo e o valor real registrado no banco).

\begin{itemize}
    \item \textbf{Cenário Real (Registrado no SINAN):} 4.993 casos de Grau 2 (9.7\% da amostra de teste).
    \item \textbf{Cenário Predito (Estimado pela IA):} 8.083 casos de Grau 2 (15.7\% da amostra de teste).
\end{itemize}

\subsection{Consequências para a Saúde Pública}
O modelo, "aprendendo" com quase 20 anos de histórico clínico, identificou que, com base nas características apresentadas pelos pacientes (idade, baciloscopia, número de nervos afetados, forma clínica), **o número de casos com incapacidade grave deveria ser aproximadamente 60\% maior do que o registrado oficialmente**.

**Interpretação Epidemiológica:** Esse fenômeno não deve ser interpretado como uma redução real da virulência da doença. Pelo contrário, aponta fortemente para uma **subnotificação qualificada** e sistemática. Durante e após a pandemia, a sobrecarga e desestruturação do sistema de saúde podem ter levado a:
\begin{enumerate}
    \item **Simplificação da Avaliação:** A Avaliação Neurológica Simplificada, necessária para classificar o GIF, exige tempo e contato físico. Em um cenário de distanciamento e sobrecarga, exames físicos detalhados podem ter sido negligenciados.
    \item **Preenchimento Padrão (Default):** Uso de valores padrão ("Grau 0" ou "Grau 1") em fichas de notificação devido à falta de tempo ou rotatividade de profissionais não treinados.
    \item **Barreiras de Acesso:** Pacientes com complicações graves podem ter tido dificuldades em acessar centros de referência para a avaliação final na alta, resultando em altas administrativas sem avaliação adequada.
\end{enumerate}

Essa "demanda reprimida" oculta representa um grave risco de saúde pública. Pacientes com incapacidades não reconhecidas nos registros deixam de ser priorizados para encaminhamento a serviços de reabilitação física e prevenção de incapacidades. Isso aumenta o risco de desenvolvimento de deformidades permanentes, exclusão social e perpetuação do estigma, além de gerar custos futuros mais elevados para o sistema de saúde \citep{oliveira2024mobilidade}.

% ==================================================================
% CAPÍTULO 5: CONCLUSÃO
% ==================================================================
\chapter{Conclusão}

Este estudo demonstrou a eficácia e a importância crítica da aplicação de inteligência artificial na vigilância em saúde, indo além da simples predição clínica para atuar como uma ferramenta de auditoria da qualidade da informação.

\begin{enumerate}
    \item A estratégia de **Seleção de Variáveis (RFE)** confirmou a primazia de indicadores clínicos diretos (dano neural, forma clínica, baciloscopia) sobre fatores demográficos na predição de sequelas na alta, orientando o foco da atenção clínica.
    \item O algoritmo **LightGBM** consolidou-se como a ferramenta de escolha para este domínio, oferecendo o melhor equilíbrio entre sensibilidade e precisão, superando a complexidade da GRU em termos de custo-benefício para dados tabulares do SINAN.
    \item O estudo forneceu evidências quantitativas robustas de uma **subnotificação massiva** de incapacidades físicas graves no Brasil pós-2020. Estima-se um déficit de mais de 3.000 diagnósticos de Grau 2 apenas na amostra analisada, o que sugere um problema de escala nacional.
\end{enumerate}

Recomenda-se a integração de modelos preditivos semelhantes ao ecossistema do SINAN. Tais modelos poderiam atuar como "sentinelas digitais", emitindo alertas automáticos para gestores municipais e estaduais sempre que houver inconsistência significativa entre o quadro clínico notificado e o grau de incapacidade atribuído, desencadeando ações imediatas de revisão de prontuários e busca ativa de pacientes para reavaliação.

\bibliographystyle{plainnat}
\bibliography{references_3}

\end{document}